\documentclass[11pt]{letter}
\usepackage[utf8]{inputenc}
\usepackage[T1]{fontenc}
\usepackage{mathptmx}
\usepackage[margin=1in]{geometry}
\usepackage{xcolor}
\usepackage[colorlinks=true,urlcolor=blue!70!black]{hyperref}

\signature{Hameem Mahdi, B.S.C.S., M.S.E., Ph.D.\\Mayo Clinic\\Rochester, MN 55905, USA\\mahdi.hameem@mayo.edu\\ORCID: 0009-0007-0005-3080}

\address{Hameem Mahdi\\Mayo Clinic\\Rochester, MN 55905, USA\\mahdi.hameem@mayo.edu}

\date{April 7, 2026}

\begin{document}

\begin{letter}{%
Editorial Office\\
\emph{Artificial Intelligence Review}\\
Springer Nature}

\opening{Dear Editors,}

I am writing to submit the manuscript entitled ``Perceive, Plan, Act, Self-Correct: An Architectural Framework for Goal-Directed Agentic AI Systems'' for consideration as a Review Article in \emph{Artificial Intelligence Review}.

\textbf{Summary.}
Large language model (LLM) agents that autonomously perceive environments, plan multi-step strategies, execute actions through external tools, and self-correct from feedback represent a paradigm shift in AI. Yet the field lacks a unified architectural vocabulary, forcing practitioners to navigate a ``framework zoo'' of incompatible abstractions. This manuscript proposes \textsc{Perceive--Plan--Act--Self-Correct} (PPAS), a four-phase canonical loop grounded in classical agent theory (BDI, SOAR, OODA) and extended for LLM-era systems. We decompose agentic architectures into an 8-layer technology stack and map 15+ open-source frameworks to a common reference model.

\textbf{Key contributions.}
\begin{enumerate}
    \item A formal four-phase agent loop (PPAS) that unifies existing design patterns under a single architectural vocabulary.
    \item An 8-layer technology stack mapping orchestration frameworks, memory systems, tool ecosystems, and inter-agent protocols to a common reference.
    \item Three empirical studies: (a)~a benchmark meta-analysis across SWE-Bench Verified, WebArena, and GAIA covering 12 frontier agents; (b)~a controlled design-pattern comparison demonstrating that reflective planning with human-in-the-loop approval achieves the highest task-completion rates (72.5\% on SWE-Bench) while reducing irreversible-action failures by 73\%; and (c)~an inter-agent protocol evaluation of MCP, A2A, and AG-UI.
    \item A failure-mode taxonomy and EU AI Act governance mapping for responsible agentic deployment.
\end{enumerate}

\textbf{Significance and fit for Artificial Intelligence Review.}
This manuscript is a comprehensive architectural framework and survey paper---the type of work that is at the core of \emph{Artificial Intelligence Review}'s editorial scope. The paper synthesises classical agent theory with modern LLM-era implementations, maps the complete technology stack across 15+ frameworks, and provides empirical validation through benchmark meta-analysis. This breadth and depth of review, combined with novel framework contributions, makes the manuscript well-suited for AIR's readership of AI researchers and practitioners seeking comprehensive reference works.

\textbf{Prior submission disclosure.}
In the interest of transparency, this manuscript was previously submitted to \emph{Nature Machine Intelligence} (Manuscript \#NATMACHINTELL-A26041846, 3~April 2026) and was desk-rejected on 7~April 2026 without external peer review. The Chief Editor, Dr Liesbeth Venema, stated that the decision ``in no way reflects any doubts about the quality of the work reported'' but that the paper's comprehensive review and framework character was better suited to another venue. No revisions have been requested or made; the manuscript content is unchanged.

\textbf{Preprint disclosure.}
A preprint of this work has been posted on engrXiv (DOI: \href{https://doi.org/10.31224/6738}{10.31224/6738}, March 2026), in accordance with Springer Nature's preprint policy.

\textbf{Data and code availability.}
All datasets, benchmark results, and analysis scripts are publicly available at:\\ \url{https://github.com/HAMEEMM/ai_systems_landscape_2026/tree/main/publications/01-agentic-ai}.

This manuscript is not under consideration by another journal. I have no competing interests to declare.

Thank you for your consideration. I look forward to your response.

\closing{Sincerely,}

\end{letter}
\end{document}
